\documentclass[11pt]{article}

\usepackage{amsmath,amssymb}
\usepackage{xurl}
\usepackage[hidelinks]{hyperref}
\setlength{\emergencystretch}{2em}

% Shared commands for notes in docs/paper-gaps/.
%
% Two halves.  The link commands below are project identity: OWNER/REPO is the
% placeholder that scripts/bootstrap_project.py replaces with this project's
% GitHub slug and Pages site.  Everything after them is NOTATION, and notation
% belongs to the source paper: the definitions here are an example set, and a
% project replaces them with the macros of its own paper's style file.  The one
% rule that never changes: a command defined here must mean exactly what the
% source paper's macro means, or a note silently says something else.

\newcommand{\Lean}{\textsc{Lean}}
\newcommand{\leanid}[1]{\textsf{#1}}
\newcommand{\ghissue}[1]{\href{https://github.com/OWNER/REPO/issues/#1}{GitHub issue \##1}}
\newcommand{\ghpr}[1]{\href{https://github.com/OWNER/REPO/pull/#1}{PR \##1}}
% Local issue tracker (issues/NNNN-*.md in this repository); no remote URL.
\newcommand{\localissue}[1]{issue \##1 (\path{issues/})}
\newcommand{\doclink}[1]{\href{https://OWNER.github.io/REPO/docs/#1.html}{\path{#1}}}

% --- Notation: replace with the source paper's own macros --------------------
% \F, \N, \C, \R, \tr, \ot, \eps, \E, \bx, \bu, \bv, \by

\newcommand{\F}{\mathbb{F}}
\newcommand{\N}{\mathbb{N}}
\newcommand{\C}{\mathbb{C}}
\newcommand{\R}{\mathbb{R}}
\newcommand{\tr}{\mathrm{tr}}
\newcommand{\eps}{\epsilon}
\newcommand{\ot}{\otimes}
\newcommand{\E}{\mathop{\bf E\/}}

% bold random variables
\newcommand{\bu}{\boldsymbol{u}}
\newcommand{\bv}{\boldsymbol{v}}
\newcommand{\bx}{{\boldsymbol{x}}}
\newcommand{\by}{\boldsymbol{y}}

% T1 so literal < and > in \texttt placeholders typeset as themselves
\usepackage[T1]{fontenc}

% bra-ket
\usepackage{braket}

% --- Example structured spaces ---
\newcommand{\polyfunc}[3]{\mathcal{P}(#1, #2, #3)}
\newcommand{\polymeas}[3]{\mathrm{PolyMeas}(#1, #2, #3)}
\newcommand{\polysub}[3]{\mathrm{PolySub}(#1, #2, #3)}

% Approximation relations (paper uses \approx_\eps and \simeq_\zeta)
\newcommand{\approxeps}{\approx_{\varepsilon}}
\newcommand{\simgeq}{\simeq_{\zeta}}


\newcommand{\Id}{\mathrm{Id}}
\renewcommand{\ghissue}[1]{%
  \href{https://github.com/Dengnifer/MIPStarRE-A/issues/#1}{GitHub issue \##1}}

\title{The Binary Factor Index in the Qudit-to-Qubit Isomorphism}
\author{}
\date{2026-09-04}

\begin{document}
\maketitle

\section{At a glance}

\begin{itemize}
  \item \textbf{Difficulty.} Low.  The printed upper index is incompatible
  with both the binary coordinate array and the target tensor power.
  \item \textbf{Estimated weight.} One statement correction and no additional
  proof beyond the existing coordinatewise isomorphism.
  \item \textbf{Mathlib/project split.} The correction is entirely
  project-local notation; the proof uses finite tensor products and the fixed
  self-dual basis.
  \item \textbf{Key mathematical inputs.} A field of size $q=2^k$ has $k$
  binary coordinates, and one $q$-dimensional qudit is identified with $k$
  qubits.  This audit is tracked in \ghissue{173}.
\end{itemize}

\section{Key theorem forms}

For $q=2^k$ and $u\in\F_q^L$, write
$u_i=\sum_{j=1}^k u_{ij}e_j$ in the fixed self-dual basis.  The intended
conjugation identity is
\[
  \tau_u^W
  =\phi^\dagger
    \left(\bigotimes_{i=1}^L\bigotimes_{j=1}^k
      \sigma^W_{u_{ij}}\right)\phi,
\]
where $\phi:(\C^q)^{\ot L}\to(\C^2)^{\ot Lk}$.  The inner tensor index must
range from $1$ through $k$, not from $1$ through $q$.

\section{The source assertion and discrepancy}

Lemma \path{lem:pauli-binary} of \cite[Section~4.6]{Ji2020MIPStar} first gives
the correct tensor product with $j\in\{1,\ldots,k\}$, and defines each
$u_{ij}$ using the $k$ basis elements $e_1,\ldots,e_k$.  Its final sentence,
however, describes the $(i,j)$-th qubit factor with
$j\in\{1,\ldots,q\}$.%
\footnote{See \path{references/qpbt-paper/04_preliminaries.tex},
lines~1172--1198.}
For $k>1$ this asks for coordinates $u_{ij}$ that were not defined and for
$Lq$ factors in a target containing only $Lk$ qubits.  The proof immediately
returns to the $k$ binary coordinates.  The occurrence of $q$ is therefore a
typographical error in the mathematical statement.

\section{The correction}

Replace the final range by $j\in\{1,\ldots,k\}$.  Then the factor numbered
$s=(i-1)k+j$ lies between $1$ and $Lk$, and every binary coordinate of every
qudit appears exactly once.  This agrees with the displayed tensor product and
with the proof by tensoring the one-qudit coordinate isomorphism.

\section{Blueprint and formalization status}

The blueprint statement of \path{lem:pauli-binary} uses the corrected range
$1\leq j\leq k$, and its proof and source-correction remark cite this note.
The formal theorem indexes the inner factors by the binary-basis dimension and
therefore matches the corrected statement.  Its proof is complete in the
current tree.%
\footnote{The linked declaration is \leanid{exists\_qubitIsometry} in
\doclink{MIPStarRE/QPBT/Algebra/PauliTheorems.lean}.}

\section{Conclusion}

The upper bound $q$ in the source's final index range is a typographical error.
The correct bound is $k=\log_2 q$, as required by the binary coordinate map,
the target dimension, and the proof.

\bibliographystyle{alpha}
\bibliography{references}

\end{document}
